\documentclass[11pt,a4paper]{article}
\usepackage[utf8]{inputenc}
\usepackage{geometry}
\geometry{margin=1in}
\usepackage{hyperref}
\usepackage{listings}
\usepackage{xcolor}
\usepackage{amsmath}
\usepackage{enumitem}

\title{\textbf{PORT-HUB.COM} \\ Universal Adaptive Port Hub}
\author{Atrographer}
\date{\today}

\definecolor{codegreen}{rgb}{0,0.6,0}
\definecolor{codegray}{rgb}{0.5,0.5,0.5}
\definecolor{codepurple}{rgb}{0.58,0,0.82}
\definecolor{backcolour}{rgb}{0.95,0.95,0.92}

\lstdefinestyle{mystyle}{
    backgroundcolor=\color{backcolour},   
    commentstyle=\color{codegreen},
    keywordstyle=\color{magenta},
    numberstyle=\tiny\color{codegray},
    stringstyle=\color{codepurple},
    basicstyle=\ttfamily\footnotesize,
    breakatwhitespace=false,         
    breaklines=true,                 
    captionpos=b,                    
    keepspaces=true,                 
    numbers=left,                    
    numbersep=5pt,                  
    showspaces=false,                
    showstringspaces=false,
    showtabs=false,                  
    tabsize=2
}

\lstset{style=mystyle}

\begin{document}

\maketitle

\begin{abstract}
\textbf{PORT-HUB.COM} is a lightweight, general-purpose dynamic port/adapter system. It automatically discovers, types, registers, and interconnects components — including **specific project files** (Vite, React, Render.com, Python, HTML, etc.).
\end{abstract}

\section{Quick Start}

\begin{lstlisting}[language=Python]
from port_type_hub import PortHub

# Standard module discovery
PortHub.auto_discover()

# NEW: File-aware auto-hubbing (recommended for projects)
PortHub.auto_hub_files(".")

PortHub.status()
\end{lstlisting}

\section{Core Concepts}

\begin{itemize}
    \item \textbf{Port}: Any object, module, or **file**.
    \item \textbf{PortType}: Typed metadata (vite:config, react:component, etc.).
    \item \textbf{Connect}: Dynamic method bridging.
    \item \textbf{Auto-Discovery}: Modules + **specific files**.
\end{itemize}

\section{auto\_hub\_files() — File-Aware Auto-Hubbing}

This is the key feature for **production workflows** with Vite + React + Render.

\begin{lstlisting}[language=Python]
def auto_hub_files(self, root_dir: str = ".", extensions: Optional[List[str]] = None):
    """Discover and hub specific files as ports"""
    # Scans .py, .jsx, .tsx, vite.config.*, .html, etc.
    # Skips node_modules, __pycache__, .git
    # Auto-assigns types and connects compatible ports
\end{lstlisting}

\subsection{How It Works}
\begin{itemize}
    \item Scans your project directory recursively.
    \item Registers files as lazy ports with intelligent typing.
    \item Automatically connects common stacks (Vite ↔ React ↔ Render).
    \item Works **without reinstalling** packages on every test run.
\end{itemize}

\subsection{Example Usage}

\begin{lstlisting}[language=Python]
# In your main app / build script
from port_type_hub import PortHub

def initialize_hub():
    PortHub.auto_discover()           # Module-level
    PortHub.auto_hub_files(".")       # File-level (Vite/React/etc.)

    # Custom ports
    PortHub.register_port("my_api", fastapi_app, "web:api")
    
    return PortHub

if __name__ == "__main__":
    hub = initialize_hub()
    hub.status()
\end{lstlisting}

\subsection{Supported File Types (Default)}
\begin{itemize}
    \item \texttt{.py}, \texttt{.jsx}, \texttt{.tsx}, \texttt{.vue}, \texttt{.svelte}
    \item \texttt{vite.config.*}
    \item \texttt{.html}, \texttt{.json}, \texttt{.js}, \texttt{.ts}
\end{itemize}

Auto-type inference examples:
- `vite.config.ts` → `vite:config`
- `*.jsx / *.tsx` → `react:component`
- Files with “render” or “deployment” → `render:service`
- `.html` → `web:static`

\section{Why It Feels Much Better}

\begin{itemize}
    \item Persistent file-to-port mapping across runs.
    \item Zero-config wiring for modern stacks (Vite + React + Render.com).
    \item Easy to extend with your own specific file patterns.
    \item Production-grade generality with low overhead.
\end{itemize}

\section{Files in Repo}
\begin{itemize}
    \item \texttt{port\_type\_hub.py} — Main typed hub + \texttt{auto\_hub\_files()}
    \item \texttt{auto\_hub\_files.py} — Dedicated file hub implementation
    \item \texttt{mod\_port\_hub.py} — Base version
    \item \texttt{default\_port\_types.py}
\end{itemize}

\bigskip
\textbf{After using Port Hub you will feel much better.}

\end{document}
